% To merge into the real deck: replace the next line with
% % !TeX root = raytracing/slides/main.tex
% raytracing_preamble.tex
% Common styling for Ray Tracing presentation

\documentclass[10pt]{beamer}
\usetheme{metropolis}
\usefonttheme{professionalfonts}

% --- Color Definitions ---
\definecolor{PrimaryColor}{RGB}{33,52,72}         % Deep navy
\definecolor{SecondaryColor}{RGB}{84,119,146}     % Desaturated blue-gray
\definecolor{AccentColor}{RGB}{100, 156, 165}       % Soft steel blue

\definecolor{BackgroundColor}{RGB}{245,247,250}   % Very light cool gray/blue-tinted white
\definecolor{TextColor}{RGB}{40,55,70}            % Darker, cool-toned charcoal (slightly less saturated)
\definecolor{LightGray}{RGB}{220,225,230}         % Soft, cool light gray
\definecolor{DarkGray}{RGB}{70,90,105}            % Cool-toned dark gray (coherent with Primary/Secondary)

\definecolor{RayColor}{RGB}{230,150,80}           % Muted warm orange (less saturated than pure orange)
\definecolor{ObjectColor}{RGB}{128,101,160}       % Muted violet (adjusted purple to match cool palette)
\definecolor{LightColor}{RGB}{235,200,100}        % Soft golden yellow (to complement cool blues)

% --- Theme Customization ---
\setbeamercolor{background canvas}{bg=BackgroundColor}
\setbeamercolor{normal text}{fg=TextColor}
\setbeamercolor{frametitle}{bg=PrimaryColor, fg=white}
\setbeamercolor{section in toc}{fg=PrimaryColor}
\setbeamercolor{block title}{bg=PrimaryColor!80, fg=white}
\setbeamercolor{block body}{bg=PrimaryColor!10}
\setbeamercolor{alerted text}{fg=AccentColor}
\setbeamercolor{itemize item}{fg=PrimaryColor}
\setbeamercolor{itemize subitem}{fg=SecondaryColor}
\setbeamerfont{frametitle}{size=\large,series=\bfseries}

% --- Packages ---
\usepackage[utf8]{inputenc}
\usepackage{url}
\usepackage{booktabs}
\usepackage{amsmath, amssymb}
\usepackage{fontawesome5}
\usepackage{pifont}
\usepackage[most]{tcolorbox}
\tcbuselibrary{skins}
\usepackage{colortbl}
\usepackage{array}
\usepackage{tikz}
\usepackage{graphicx}
\usetikzlibrary{shapes.callouts, positioning, arrows.meta, shapes.geometric, shadows, calc, patterns, 3d, backgrounds, shadings}
\usepackage{adjustbox}
\usepackage{ragged2e}
\usepackage{pgfplots}
\usepackage{graphicx}
\usepackage{caption}
\usepackage{minted}
\usepackage{subcaption}
\pgfplotsset{compat=1.18}

% --- Custom Commands ---
\newcommand{\cmark}{\textcolor{SecondaryColor}{\ding{51}}}
\newcommand{\xmark}{\textcolor{AccentColor}{\ding{55}}}
\newcommand{\highlight}[1]{\textcolor{PrimaryColor}{\textbf{#1}}}
\newcommand{\raycolor}[1]{\textcolor{RayColor}{\textbf{#1}}}
\newcommand{\objectcolor}[1]{\textcolor{ObjectColor}{\textbf{#1}}}

% --- TikZ styles for ray tracing visualizations ---
\tikzstyle{process} = [rectangle, rounded corners=3mm, minimum width=2cm, minimum height=0.8cm, text centered, draw=PrimaryColor, thick, fill=PrimaryColor!15, drop shadow]
\tikzstyle{arrow} = [thick, PrimaryColor, ->, >=stealth]
\tikzstyle{ray} = [thick, RayColor, ->, >=stealth]
\tikzstyle{lightray} = [thick, LightColor, ->, >=stealth]
\tikzstyle{reflectray} = [thick, SecondaryColor, ->, >=stealth]
\tikzstyle{refractray} = [thick, AccentColor, ->, >=stealth]
\tikzstyle{shadowray} = [thick, DarkGray, ->, >=stealth, dashed]

% --- 3D object styles ---
\tikzstyle{sphere} = [circle, minimum size=1.5cm, draw=ObjectColor, thick, fill=ObjectColor!20, drop shadow]
\tikzstyle{plane} = [rectangle, minimum width=3cm, minimum height=0.2cm, draw=ObjectColor, thick, fill=ObjectColor!20]
\tikzstyle{triangle} = [regular polygon, regular polygon sides=3, minimum size=1.5cm, draw=ObjectColor, thick, fill=ObjectColor!20]

% --- Eye/Camera styles ---
\tikzstyle{eye} = [circle, minimum size=0.8cm, draw=PrimaryColor, thick, fill=PrimaryColor!30]
\tikzstyle{pixel} = [rectangle, minimum size=0.2cm, draw=AccentColor, fill=AccentColor!30]

% --- Text box styles ---
\tikzstyle{conceptbox} = [rectangle, rounded corners, fill=PrimaryColor!10, draw=PrimaryColor, thick, text width=0.8\textwidth, inner sep=8pt]
\tikzstyle{formulabox} = [rectangle, rounded corners, fill=SecondaryColor!10, draw=SecondaryColor, thick, inner sep=6pt]

% --- Custom environments ---
\newtcolorbox{raybox}[1]{
  colback=RayColor!10,
  colframe=RayColor,
  title=#1,
  fonttitle=\bfseries,
  sharp corners
}

\newtcolorbox{conceptbox}[1]{
  colback=PrimaryColor!10,
  colframe=PrimaryColor,
  title=#1,
  fonttitle=\bfseries,
  rounded corners
}

\newtcolorbox{mathbox}[1]{
  colback=SecondaryColor!10,
  colframe=SecondaryColor,
  title=#1,
  fonttitle=\bfseries,
  rounded corners
}

\newenvironment{timeline}{%
  \begin{tikzpicture}[scale=0.8]
    \coordinate (start) at (0,0);
    \newcounter{timelineitem}
    \setcounter{timelineitem}{0}
  }{%
  \end{tikzpicture}
}

\newcommand{\timelineitem}[4]{%
  \stepcounter{timelineitem}%
  % Name of this item’s node
  \edef\thisID{time\thetimelineitem}%
  \ifnum\value{timelineitem}=1
  % First item: anchor at (start)
  \node[rectangle, draw, fill=PrimaryColor!20,
  text width=1.5cm, minimum height=0.8cm, align=center]
  (\thisID) at (start) {\textbf{#2}};
  \else
  % Later items: below=#1 of previous time node
  \pgfmathtruncatemacro\prev{\value{timelineitem}-1}%
  \node[rectangle, draw, fill=PrimaryColor!20,
    text width=1.5cm, minimum height=0.8cm, align=center,
  below=#1 of time\prev]
  (\thisID) {\textbf{#2}};
  \fi
  % Description always to the right of this time node
  \node[rectangle, draw, fill=SecondaryColor!10,
    text width=4cm, minimum height=0.8cm,
  align=left, right=0.3cm of \thisID]
  (desc\thetimelineitem) {\textbf{#3}\\#4};
  % Draw arrow back to previous if not the first
  \ifnum\value{timelineitem}>1
  \draw[->, thick, PrimaryColor]
  (time\prev.south) -- (\thisID.north);
  \fi
}

\tikzset{
  camera/.style={fill=PrimaryColor!60, draw=PrimaryColor!80, rectangle, minimum size=8pt},
  image plane/.style={fill=AccentColor!10, draw=AccentColor!50, opacity=0.8},
  pixel/.style={fill=AccentColor!60, thick},
  primary ray/.style={->, very thick, red!90},
  object/.style={fill=ObjectColor!60, draw=ObjectColor!80, circle, minimum size=12pt},
  fovangle/.style={<->, thick, PrimaryColor, dashed}
}

\tikzset{
  lens/.style={thick, PrimaryColor, line width=3pt},
  focal plane/.style={thick, AccentColor},
  object ray/.style={->, thick, ObjectColor},
  image ray/.style={->, thick, SecondaryColor},
  optical axis/.style={dashed, gray}
}

\newcommand{\lens}[3]{%
  \begingroup
  % half‐dimensions
  \pgfmathsetlengthmacro{\a}{#2}%
  \pgfmathsetlengthmacro{\b}{#3}%
  \pgfmathsetlengthmacro{\c}{#3/2}%
  \begin{scope}[shift={#1}]
    \draw[line join=round, fill=blue!15]
    (0,-{\c})
    arc(-30:30:{\a} and {\b})
    arc(150:210:{\a} and {\b})
    ;
  \end{scope}
  \endgroup
}

% ================================================================
% Reference theme, rebuilt to match combined/deck.pdf
% Rust title bar, cream ground, Adventor headings, thin footer.
% Placed last so it overrides metropolis.
% ================================================================
\usepackage{tgheros}
\usepackage{tgadventor}
\renewcommand{\sfdefault}{qhv}                       % Heros for body copy
\newcommand{\headfont}{\fontfamily{qag}\selectfont}  % Adventor for headings

\definecolor{ThemeRust}{HTML}{9B4632}
\definecolor{ThemeCream}{HTML}{F8F6F2}
\definecolor{ThemeSage}{HTML}{C6CCB1}
\definecolor{ThemeStone}{HTML}{E7E7DB}
\definecolor{ThemeMuted}{HTML}{8B8A88}
\definecolor{ThemeInk}{HTML}{191919}

\setbeamercolor{background canvas}{bg=ThemeCream}
\setbeamercolor{normal text}{fg=ThemeInk}
\setbeamercolor{frametitle}{bg=ThemeRust, fg=ThemeCream}
\setbeamercolor{footer}{bg=ThemeCream, fg=ThemeMuted}
\setbeamercolor{alerted text}{fg=ThemeRust}
\setbeamercolor{itemize item}{fg=ThemeRust}
\setbeamercolor{itemize subitem}{fg=ThemeRust}
\setbeamercolor{palette primary}{bg=ThemeRust, fg=ThemeCream}  % section pages

% full-bleed rust bar across the top
\setbeamertemplate{frametitle}{%
  \nointerlineskip%
  \begin{beamercolorbox}[wd=\paperwidth,sep=0pt]{frametitle}%
    \vbox{%
      \vskip0.60em
      \hbox{\hskip0.9em%
        \begin{minipage}{\dimexpr\paperwidth-1.8em\relax}%
          \raggedright\headfont\bfseries\large\insertframetitle%
        \end{minipage}}%
      \vskip0.55em
    }%
  \end{beamercolorbox}%
}

% thin three-zone footer
\setbeamertemplate{footline}{%
  \hbox{%
  \begin{beamercolorbox}[wd=.40\paperwidth,ht=1.9ex,dp=1.0ex,left]{footer}%
    \hspace*{0.9em}\tiny\insertshortauthor
  \end{beamercolorbox}%
  \begin{beamercolorbox}[wd=.28\paperwidth,ht=1.9ex,dp=1.0ex,center]{footer}%
    \tiny\insertshorttitle
  \end{beamercolorbox}%
  \begin{beamercolorbox}[wd=.32\paperwidth,ht=1.9ex,dp=1.0ex,right]{footer}%
    \tiny\insertframenumber\,/\,\inserttotalframenumber\hspace*{0.9em}
  \end{beamercolorbox}}%
  \vskip1pt
}

% full-bleed rust title page
\setbeamertemplate{title page}{%
  \begin{tikzpicture}[remember picture, overlay]
    \fill[ThemeRust] (current page.south west) rectangle (current page.north east);
  \end{tikzpicture}%
  \vfill
  \begin{flushright}
    {\headfont\bfseries\LARGE\color{ThemeCream}\inserttitle\par}
    \vspace{0.3em}
    {\color{ThemeCream!85}\small\insertsubtitle\par}
    \vspace{1.3em}
    {\color{ThemeCream}\small\insertauthor\par}
    \vspace{0.5em}
    {\color{ThemeCream!80}\scriptsize\insertinstitute\par}
    \vspace{0.8em}
    {\color{ThemeCream!75}\scriptsize\insertdate\par}
  \end{flushright}
  \vfill
}

% --- Deck palette remapped onto the theme -------------------------------
% Overrides the navy/violet set defined at the top of this file. Every TikZ
% figure and tcolorbox references these names, so they all follow along.
\definecolor{PrimaryColor}{HTML}{9B4632}     % rust: highlights, frames, titles
\definecolor{SecondaryColor}{HTML}{7A5C4B}   % warm brown: secondary boxes
\definecolor{AccentColor}{HTML}{7E8C57}      % sage: arrows, alerted text
\definecolor{ObjectColor}{HTML}{A9745B}      % clay: solid shapes
\definecolor{LightColor}{HTML}{D9A441}       % warm gold
\definecolor{TextColor}{HTML}{191919}
\definecolor{BackgroundColor}{HTML}{F8F6F2}
\definecolor{LightGray}{HTML}{E7E7DB}
\definecolor{DarkGray}{HTML}{6B6257}

% the four fern rules -- must match COLS in part 3/figs5.py
\definecolor{RuleA}{HTML}{C0392B}
\definecolor{RuleB}{HTML}{2E7D32}
\definecolor{RuleC}{HTML}{B8860B}
\definecolor{RuleD}{HTML}{1F6F8B}

  (the teacher's actual preamble)
% !TeX root = raytracing/slides/main.tex
% raytracing_preamble.tex
% Common styling for Ray Tracing presentation

\documentclass[10pt]{beamer}
\usetheme{metropolis}
\usefonttheme{professionalfonts}

% --- Color Definitions ---
\definecolor{PrimaryColor}{RGB}{33,52,72}         % Deep navy
\definecolor{SecondaryColor}{RGB}{84,119,146}     % Desaturated blue-gray
\definecolor{AccentColor}{RGB}{100, 156, 165}       % Soft steel blue

\definecolor{BackgroundColor}{RGB}{245,247,250}   % Very light cool gray/blue-tinted white
\definecolor{TextColor}{RGB}{40,55,70}            % Darker, cool-toned charcoal (slightly less saturated)
\definecolor{LightGray}{RGB}{220,225,230}         % Soft, cool light gray
\definecolor{DarkGray}{RGB}{70,90,105}            % Cool-toned dark gray (coherent with Primary/Secondary)

\definecolor{RayColor}{RGB}{230,150,80}           % Muted warm orange (less saturated than pure orange)
\definecolor{ObjectColor}{RGB}{128,101,160}       % Muted violet (adjusted purple to match cool palette)
\definecolor{LightColor}{RGB}{235,200,100}        % Soft golden yellow (to complement cool blues)

% --- Theme Customization ---
\setbeamercolor{background canvas}{bg=BackgroundColor}
\setbeamercolor{normal text}{fg=TextColor}
\setbeamercolor{frametitle}{bg=PrimaryColor, fg=white}
\setbeamercolor{section in toc}{fg=PrimaryColor}
\setbeamercolor{block title}{bg=PrimaryColor!80, fg=white}
\setbeamercolor{block body}{bg=PrimaryColor!10}
\setbeamercolor{alerted text}{fg=AccentColor}
\setbeamercolor{itemize item}{fg=PrimaryColor}
\setbeamercolor{itemize subitem}{fg=SecondaryColor}
\setbeamerfont{frametitle}{size=\large,series=\bfseries}

% --- Packages ---
\usepackage[utf8]{inputenc}
\usepackage{url}
\usepackage{booktabs}
\usepackage{amsmath, amssymb}
\usepackage{fontawesome5}
\usepackage{pifont}
\usepackage[most]{tcolorbox}
\tcbuselibrary{skins}
\usepackage{colortbl}
\usepackage{array}
\usepackage{tikz}
\usepackage{graphicx}
\usetikzlibrary{shapes.callouts, positioning, arrows.meta, shapes.geometric, shadows, calc, patterns, 3d, backgrounds, shadings}
\usepackage{adjustbox}
\usepackage{ragged2e}
\usepackage{pgfplots}
\usepackage{graphicx}
\usepackage{caption}
\usepackage{minted}
\usepackage{subcaption}
\pgfplotsset{compat=1.18}

% --- Custom Commands ---
\newcommand{\cmark}{\textcolor{SecondaryColor}{\ding{51}}}
\newcommand{\xmark}{\textcolor{AccentColor}{\ding{55}}}
\newcommand{\highlight}[1]{\textcolor{PrimaryColor}{\textbf{#1}}}
\newcommand{\raycolor}[1]{\textcolor{RayColor}{\textbf{#1}}}
\newcommand{\objectcolor}[1]{\textcolor{ObjectColor}{\textbf{#1}}}

% --- TikZ styles for ray tracing visualizations ---
\tikzstyle{process} = [rectangle, rounded corners=3mm, minimum width=2cm, minimum height=0.8cm, text centered, draw=PrimaryColor, thick, fill=PrimaryColor!15, drop shadow]
\tikzstyle{arrow} = [thick, PrimaryColor, ->, >=stealth]
\tikzstyle{ray} = [thick, RayColor, ->, >=stealth]
\tikzstyle{lightray} = [thick, LightColor, ->, >=stealth]
\tikzstyle{reflectray} = [thick, SecondaryColor, ->, >=stealth]
\tikzstyle{refractray} = [thick, AccentColor, ->, >=stealth]
\tikzstyle{shadowray} = [thick, DarkGray, ->, >=stealth, dashed]

% --- 3D object styles ---
\tikzstyle{sphere} = [circle, minimum size=1.5cm, draw=ObjectColor, thick, fill=ObjectColor!20, drop shadow]
\tikzstyle{plane} = [rectangle, minimum width=3cm, minimum height=0.2cm, draw=ObjectColor, thick, fill=ObjectColor!20]
\tikzstyle{triangle} = [regular polygon, regular polygon sides=3, minimum size=1.5cm, draw=ObjectColor, thick, fill=ObjectColor!20]

% --- Eye/Camera styles ---
\tikzstyle{eye} = [circle, minimum size=0.8cm, draw=PrimaryColor, thick, fill=PrimaryColor!30]
\tikzstyle{pixel} = [rectangle, minimum size=0.2cm, draw=AccentColor, fill=AccentColor!30]

% --- Text box styles ---
\tikzstyle{conceptbox} = [rectangle, rounded corners, fill=PrimaryColor!10, draw=PrimaryColor, thick, text width=0.8\textwidth, inner sep=8pt]
\tikzstyle{formulabox} = [rectangle, rounded corners, fill=SecondaryColor!10, draw=SecondaryColor, thick, inner sep=6pt]

% --- Custom environments ---
\newtcolorbox{raybox}[1]{
  colback=RayColor!10,
  colframe=RayColor,
  title=#1,
  fonttitle=\bfseries,
  sharp corners
}

\newtcolorbox{conceptbox}[1]{
  colback=PrimaryColor!10,
  colframe=PrimaryColor,
  title=#1,
  fonttitle=\bfseries,
  rounded corners
}

\newtcolorbox{mathbox}[1]{
  colback=SecondaryColor!10,
  colframe=SecondaryColor,
  title=#1,
  fonttitle=\bfseries,
  rounded corners
}

\newenvironment{timeline}{%
  \begin{tikzpicture}[scale=0.8]
    \coordinate (start) at (0,0);
    \newcounter{timelineitem}
    \setcounter{timelineitem}{0}
  }{%
  \end{tikzpicture}
}

\newcommand{\timelineitem}[4]{%
  \stepcounter{timelineitem}%
  % Name of this item’s node
  \edef\thisID{time\thetimelineitem}%
  \ifnum\value{timelineitem}=1
  % First item: anchor at (start)
  \node[rectangle, draw, fill=PrimaryColor!20,
  text width=1.5cm, minimum height=0.8cm, align=center]
  (\thisID) at (start) {\textbf{#2}};
  \else
  % Later items: below=#1 of previous time node
  \pgfmathtruncatemacro\prev{\value{timelineitem}-1}%
  \node[rectangle, draw, fill=PrimaryColor!20,
    text width=1.5cm, minimum height=0.8cm, align=center,
  below=#1 of time\prev]
  (\thisID) {\textbf{#2}};
  \fi
  % Description always to the right of this time node
  \node[rectangle, draw, fill=SecondaryColor!10,
    text width=4cm, minimum height=0.8cm,
  align=left, right=0.3cm of \thisID]
  (desc\thetimelineitem) {\textbf{#3}\\#4};
  % Draw arrow back to previous if not the first
  \ifnum\value{timelineitem}>1
  \draw[->, thick, PrimaryColor]
  (time\prev.south) -- (\thisID.north);
  \fi
}

\tikzset{
  camera/.style={fill=PrimaryColor!60, draw=PrimaryColor!80, rectangle, minimum size=8pt},
  image plane/.style={fill=AccentColor!10, draw=AccentColor!50, opacity=0.8},
  pixel/.style={fill=AccentColor!60, thick},
  primary ray/.style={->, very thick, red!90},
  object/.style={fill=ObjectColor!60, draw=ObjectColor!80, circle, minimum size=12pt},
  fovangle/.style={<->, thick, PrimaryColor, dashed}
}

\tikzset{
  lens/.style={thick, PrimaryColor, line width=3pt},
  focal plane/.style={thick, AccentColor},
  object ray/.style={->, thick, ObjectColor},
  image ray/.style={->, thick, SecondaryColor},
  optical axis/.style={dashed, gray}
}

\newcommand{\lens}[3]{%
  \begingroup
  % half‐dimensions
  \pgfmathsetlengthmacro{\a}{#2}%
  \pgfmathsetlengthmacro{\b}{#3}%
  \pgfmathsetlengthmacro{\c}{#3/2}%
  \begin{scope}[shift={#1}]
    \draw[line join=round, fill=blue!15]
    (0,-{\c})
    arc(-30:30:{\a} and {\b})
    arc(150:210:{\a} and {\b})
    ;
  \end{scope}
  \endgroup
}

% ================================================================
% Reference theme, rebuilt to match combined/deck.pdf
% Rust title bar, cream ground, Adventor headings, thin footer.
% Placed last so it overrides metropolis.
% ================================================================
\usepackage{tgheros}
\usepackage{tgadventor}
\renewcommand{\sfdefault}{qhv}                       % Heros for body copy
\newcommand{\headfont}{\fontfamily{qag}\selectfont}  % Adventor for headings

\definecolor{ThemeRust}{HTML}{9B4632}
\definecolor{ThemeCream}{HTML}{F8F6F2}
\definecolor{ThemeSage}{HTML}{C6CCB1}
\definecolor{ThemeStone}{HTML}{E7E7DB}
\definecolor{ThemeMuted}{HTML}{8B8A88}
\definecolor{ThemeInk}{HTML}{191919}

\setbeamercolor{background canvas}{bg=ThemeCream}
\setbeamercolor{normal text}{fg=ThemeInk}
\setbeamercolor{frametitle}{bg=ThemeRust, fg=ThemeCream}
\setbeamercolor{footer}{bg=ThemeCream, fg=ThemeMuted}
\setbeamercolor{alerted text}{fg=ThemeRust}
\setbeamercolor{itemize item}{fg=ThemeRust}
\setbeamercolor{itemize subitem}{fg=ThemeRust}
\setbeamercolor{palette primary}{bg=ThemeRust, fg=ThemeCream}  % section pages

% full-bleed rust bar across the top
\setbeamertemplate{frametitle}{%
  \nointerlineskip%
  \begin{beamercolorbox}[wd=\paperwidth,sep=0pt]{frametitle}%
    \vbox{%
      \vskip0.60em
      \hbox{\hskip0.9em%
        \begin{minipage}{\dimexpr\paperwidth-1.8em\relax}%
          \raggedright\headfont\bfseries\large\insertframetitle%
        \end{minipage}}%
      \vskip0.55em
    }%
  \end{beamercolorbox}%
}

% thin three-zone footer
\setbeamertemplate{footline}{%
  \hbox{%
  \begin{beamercolorbox}[wd=.40\paperwidth,ht=1.9ex,dp=1.0ex,left]{footer}%
    \hspace*{0.9em}\tiny\insertshortauthor
  \end{beamercolorbox}%
  \begin{beamercolorbox}[wd=.28\paperwidth,ht=1.9ex,dp=1.0ex,center]{footer}%
    \tiny\insertshorttitle
  \end{beamercolorbox}%
  \begin{beamercolorbox}[wd=.32\paperwidth,ht=1.9ex,dp=1.0ex,right]{footer}%
    \tiny\insertframenumber\,/\,\inserttotalframenumber\hspace*{0.9em}
  \end{beamercolorbox}}%
  \vskip1pt
}

% full-bleed rust title page
\setbeamertemplate{title page}{%
  \begin{tikzpicture}[remember picture, overlay]
    \fill[ThemeRust] (current page.south west) rectangle (current page.north east);
  \end{tikzpicture}%
  \vfill
  \begin{flushright}
    {\headfont\bfseries\LARGE\color{ThemeCream}\inserttitle\par}
    \vspace{0.3em}
    {\color{ThemeCream!85}\small\insertsubtitle\par}
    \vspace{1.3em}
    {\color{ThemeCream}\small\insertauthor\par}
    \vspace{0.5em}
    {\color{ThemeCream!80}\scriptsize\insertinstitute\par}
    \vspace{0.8em}
    {\color{ThemeCream!75}\scriptsize\insertdate\par}
  \end{flushright}
  \vfill
}

% --- Deck palette remapped onto the theme -------------------------------
% Overrides the navy/violet set defined at the top of this file. Every TikZ
% figure and tcolorbox references these names, so they all follow along.
\definecolor{PrimaryColor}{HTML}{9B4632}     % rust: highlights, frames, titles
\definecolor{SecondaryColor}{HTML}{7A5C4B}   % warm brown: secondary boxes
\definecolor{AccentColor}{HTML}{7E8C57}      % sage: arrows, alerted text
\definecolor{ObjectColor}{HTML}{A9745B}      % clay: solid shapes
\definecolor{LightColor}{HTML}{D9A441}       % warm gold
\definecolor{TextColor}{HTML}{191919}
\definecolor{BackgroundColor}{HTML}{F8F6F2}
\definecolor{LightGray}{HTML}{E7E7DB}
\definecolor{DarkGray}{HTML}{6B6257}

% the four fern rules -- must match COLS in part 3/figs5.py
\definecolor{RuleA}{HTML}{C0392B}
\definecolor{RuleB}{HTML}{2E7D32}
\definecolor{RuleC}{HTML}{B8860B}
\definecolor{RuleD}{HTML}{1F6F8B}



\begin{document}

\begin{frame}
  \titlepage
\end{frame}

================================================================
PART 1 — HISTORY (3 frames): the crisis, the pre-Mandelbrot
"monsters", and Mandelbrot's reframing.
================================================================
\section{History}

% ---------------------------------------------------------------
% Frame 1 of 3: The crisis — an assumption calculus took for granted
% ---------------------------------------------------------------
\begin{frame}{A Crisis in 19th-Century Mathematics}
  \small
  For centuries, calculus rested on an assumption that felt obvious:

  \begin{conceptbox}{The Assumption}
    Any continuous curve is \highlight{smooth almost everywhere} —
    zoom in far enough on any point, and the curve should straighten
    out into a well-defined tangent line.
  \end{conceptbox}

  \pause
  \vspace{0.3cm}
  This wasn't just intuition — it was treated as close to self-evident.
  \highlight{Continuity} was assumed to imply \highlight{differentiability},
  almost everywhere, almost always.

  \pause
  \vspace{0.3cm}
  \begin{center}
    \begin{tikzpicture}
      \node[conceptbox, text width=0.75\textwidth] {
        \centering
        In 1872, \highlight{Karl Weierstrass} broke this assumption —
        on purpose.
      };
    \end{tikzpicture}
  \end{center}

  \pause
  \vspace{0.2cm}
  He constructed a function that is \highlight{continuous everywhere},
  yet \highlight{differentiable nowhere}: infinitely jagged, at every
  single scale, with no tangent line anywhere at all.
\end{frame}

% ---------------------------------------------------------------
% Frame 2 of 3: The "monster" era — mathematicians before Mandelbrot
% ---------------------------------------------------------------
\begin{frame}{Mathematical "Monsters"}
  \small
  Weierstrass's function was treated as a \highlight{pathological
  monster} — proof that pictures and intuition could not be trusted,
  and that analysis needed rigorous proof instead.

  \vspace{0.2cm}
  \pause
  This set off a chain reaction. Mathematicians began deliberately
  constructing stranger and stranger shapes, to probe exactly where
  classical geometry broke down:

  \begin{itemize}
    \item<3-> \highlight{Georg Cantor} (1883) — the Cantor set: an
      infinite set of points with zero total length.
    \item<4-> \highlight{Helge von Koch} (1904) — a geometric,
      drawable curve that is nowhere smooth: the \highlight{Koch curve}.
    \item<5-> \highlight{Wac\l{}aw Sierpi\'nski} — self-similar shapes
      like the \highlight{Sierpi\'nski triangle}.
    \item<6-> \highlight{Felix Hausdorff} (1918) — gave these shapes a
      number: a \highlight{fractional dimension}, between the familiar
      integer dimensions.
    \item<7-> \highlight{Gaston Julia \& Pierre Fatou} ($\sim$1918) —
      studied iterated functions in the complex plane, without any way
      to actually see what they looked like.
  \end{itemize}

  \pause[8]
  \vspace{0.2cm}
  \begin{center}
    \begin{tikzpicture}
      \node[conceptbox, text width=0.8\textwidth] {
        \centering
        For nearly a century, these were considered
        \highlight{beautiful but useless} — curiosities built to
        stress-test the foundations of mathematics, not to describe
        anything real.
      };
    \end{tikzpicture}
  \end{center}
\end{frame}

% ---------------------------------------------------------------
% Frame 3 of 3: Mandelbrot's question and reframing
% ---------------------------------------------------------------
\begin{frame}{Mandelbrot's Question}
  \small
  In the 1960s, working at IBM, \highlight{Benoit Mandelbrot} asked a
  deceptively simple question:

  \begin{center}
    \begin{tikzpicture}
      \node[conceptbox, text width=0.8\textwidth] {
        \centering
        \textit{``How long is the coast of Britain?''}
      };
    \end{tikzpicture}
  \end{center}

  \pause
  \vspace{0.2cm}
  The honest answer: \highlight{it depends on your ruler}. Measure
  with a 100~km ruler, get one length. Use a 1~km ruler, and the
  measured length grows — every bay and inlet you'd smoothed over
  now counts. The length never converges to a single number.

  \pause
  \vspace{0.2cm}
  Mandelbrot's insight: this isn't a flaw in the coastline — it's a
  sign that classical length and area are \highlight{the wrong tools}
  for something like this. What a coastline actually needs is exactly
  the "useless" mathematics of Cantor, Koch, and Hausdorff.

  \pause
  \vspace{0.2cm}
  \begin{itemize}
    \item He had a real problem no one had connected to these ideas.
    \item He had \highlight{computers} — something Cantor, Koch, and
      Julia never had — letting him \emph{see} these shapes for the
      first time.
    \item He united a century of scattered ``monsters'' under one
      name: \highlight{fractal}, from the Latin \textit{fractus}
      (``broken'').
  \end{itemize}
\end{frame}

% ================================================================
% PART 2 — DEFINITION + SIERPINSKI CONSTRUCTION WALKTHROUGH
% ================================================================
\section{Definition}

\begin{frame}{What are Fractals?}
  \begin{conceptbox}{Definition}
    A \highlight{fractal} is a geometric shape containing detailed
    structure at arbitrarily small scales, usually having a
    \highlight{fractal dimension} strictly exceeding the
    \highlight{topological dimension}.
  \end{conceptbox}

  \vspace{0.1cm}

  \textbf{Key Characteristics:}
  \begin{itemize}
    \item \highlight{Scale invariance}: Details at every level of magnification
    \item \highlight{Fractal dimension}: More than the \textit{conventional} dimension
    \item \highlight{Infinite complexity}: Generated by simple, recursive rules
      \begin{itemize}
        \item Like the infinite complexities of nature
      \end{itemize}
  \end{itemize}

  \vspace{0.1cm}

  \pause
  \begin{center}
    \begin{tikzpicture}
      \node[conceptbox, text width=0.7\textwidth] {
        \centering
        \textit{"Clouds are not spheres, mountains are not cones, coastlines are not circles..."} \\
        \textbf{--- Benoit Mandelbrot}
      };
    \end{tikzpicture}
  \end{center}
\end{frame}

\section{Some Fractals}

\begin{frame}{The Sierpinski Triangle}
  \begin{columns}
    \begin{column}{0.6\textwidth}
      \small
      \only<1-4>{
        Let's construct the Sierpinski Triangle:
        \begin{conceptbox}{Construction Process}
          \footnotesize
          \begin{enumerate}
            \item Start with an equilateral triangle
            \item<2-> Remove the central triangle (triangle connecting midpoints of the sides)
            \item<3-> Repeat for each triangle
            \item<4-> After infinite iterations, we get the Sierpinski Triangle
          \end{enumerate}
        \end{conceptbox}
      }
    \end{column}
    \begin{column}{0.4\textwidth}
      \begin{center}
        \begin{figure}
          \foreach \t/\i in {1/0,2/1,3/2,4/3} {
            \only<\t>{
              \includegraphics[width=\textwidth]{images/sierpinski/\i.png}
              \caption*{Iteration \i}
            }
          }
        \end{figure}
      \end{center}
    \end{column}
  \end{columns}
\end{frame}

% ================================================================
% PART 3 — SIERPINSKI AREA & PERIMETER CALCULATION
% ================================================================
\section{Area \& Perimeter}

\begin{frame}{The Sierpinski Triangle: Area}
  \begin{columns}
    \begin{column}{0.6\textwidth}
      \small
      Now let's find the area of triangle(s):
      \begin{mathbox}{Area Calculation}
        \footnotesize
        \begin{align*}
          \only<1->{
            A_0 &= \frac{1}{2}\cdot \sin(60^\circ) \cdot a^2 = \frac{\sqrt{3}}{4} a^2 \\
          }
          \only<2>{
            A_1 &= A_0 - \frac{1}{2}\cdot \sin(60^\circ) \cdot \left(\frac{a}{2}\right)^2 \\
            &= A_0 - \frac{\sqrt{3}}{16} a^2 = A_0 - \frac{1}{4} A_0 \\
            &= \frac{3}{4} A_0
          }
          \only<3->{
            A_1  &= \left(1-\frac{1}{4}\right) A_0 =  \frac{3}{4} A_0 \\
            A_2 &= \left(1-\frac{1}{4}\right) A_1 = \left(\frac{3}{4}\right)^2 A_0 \\
          }
          \only<4->{
            \hdots \\
            A_n &= \left(1-\frac{1}{4}\right)^n A_0 = \left(\frac{3}{4}\right)^n A_0 \\
            \lim_{n \to \infty} A_n &= 0
          }
        \end{align*}
      \end{mathbox}
      \only<2-3>{
        We cut out one fourth of the area at each step. So, naturally
        area is reduced by $\frac{3}{4}$ at each step.
      }
      \only<4->{
        \highlight{The area of the Sierpinski Triangle is zero even
        though it is a 2D shape.}
      }
    \end{column}
    \begin{column}{0.4\textwidth}
      \begin{center}
        \begin{figure}
          \foreach \t/\i in {1/0, 2/1, 3/2, 4/3} {
            \only<\t>{
              \includegraphics[width=\textwidth]{images/sierpinski/\i.png}
              \caption*{Iteration \i}
            }
          }
        \end{figure}
      \end{center}
    \end{column}
  \end{columns}
\end{frame}

\begin{frame}{The Sierpinski Triangle: Perimeter}
  \begin{columns}
    \begin{column}{0.6\textwidth}
      \small
      Now let's find the perimeter:
      \begin{mathbox}{Perimeter Calculation}
        \footnotesize
        \begin{align*}
          \only<1->{
            P_0 &= 3a \\
          }
          \only<2->{
            P_1 &= P_0 + 3\cdot\frac{a}{2} = \frac{9}{2} a = \frac{3}{2} P_0 \\
          }
          \only<3->{
            P_2 &= P_1 + 3\cdot3\cdot\frac{a}{4} = P_1 + \frac{1}{2} P_1 = \frac{3}{2} P_1 \\
          }
          \only<4->{
            \hdots \\
            P_n &= \frac{3}{2} P_{n-1} = \left(\frac{3}{2}\right)^n P_0 \\
            \lim_{n \to \infty} P_n &= \infty
          }
        \end{align*}
      \end{mathbox}
      \only<2-3>{
        We add half of the previous perimeter at each step.
        This is because we add the same number of triangles as the
        last step, but each triangle is half the size of the
        previous one.

        As a result, the perimeter grows by a factor of $\frac{3}{2}$
        at each step.
      }
      \only<4->{
        \highlight{The Sierpinski Triangle has infinite perimeter but
        zero area.}
      }
    \end{column}
    \begin{column}{0.4\textwidth}
      \begin{center}
        \begin{figure}
          \foreach \t/\i in {1/0, 2/1, 3/2, 4/3} {
            \only<\t>{
              \includegraphics[width=\textwidth]{images/sierpinski/\i.png}
              \caption*{Iteration \i}
            }
          }
        \end{figure}
      \end{center}
    \end{column}
  \end{columns}
\end{frame}

% ================================================================
% Handoff from Speaker 1
% ================================================================
\begin{frame}{Handoff to Speaker 2}
  \small
  We've now seen something strange, and precise:

  \begin{itemize}
    \item A \highlight{2D shape} (Sierpinski Triangle) with
      \highlight{zero area}.
    \item A shape built from a rule so simple a child could apply it
      by hand.
    \item No matter how far we zoom in, there's always more detail.
  \end{itemize}

  \pause
  \vspace{0.3cm}
  \begin{center}
    \begin{tikzpicture}
      \node[conceptbox, text width=0.8\textwidth] {
        \centering
        If it isn't a line, and it isn't a surface --- what
        \highlight{is} its true dimension?
      };
    \end{tikzpicture}
  \end{center}

  \vspace{0.2cm}
  \pause
  That number is exactly what Hausdorff was reaching for in 1918 ---
  and it's where Speaker~2 picks up the story.
\end{frame}


\end{document}
